\section{RACF Method}
\subsection{Calibration and robust action correction}
For an independently collected transition $Z=(s,a,s^+)$, define the fixed
nominal-model score $S(Z)=\|s^+-f_0(s,a)\|_2$. Given $n$ calibration
scores, let
\begin{equation}
 k=\lceil(n+1)(1-\alpha)\rceil,\qquad q_\alpha=S_{(k)},
 \label{eq:quantile}
\end{equation}
where $S_{(k)}$ is the order statistic and $q_\alpha=+\infty$ if $k>n$.
Conformal risk control offers a broader loss construction
\cite{angelopoulos2024crc}; a residual norm makes our
margin-to-control conversion explicit.

At each step, RACF constructs the feasible set
\begin{equation}
 \mathcal F_t(m_t)=\{a\in\mathcal A:
       h(f_\theta(s_t,a))\geq m_t,\;\theta\in\Theta_t\}
 \label{eq:feasible}
\end{equation}
and computes
\begin{equation}
 a_t=\arg\min_{a\in\mathcal F_t(m_t)}\|a-a_t^\pi\|_2^2.
 \label{eq:projection}
\end{equation}
The candidate set is a local protocol object, generated before evaluation;
runtime control reads no external dataset or service. In the main study,
$\Theta=\{1500,1800\}\!\times\!\{0,3\}\!\times\!\{0.7,1.0\}$
contains eight mass (kg), grade (degrees), and actuator-gain hypotheses.
The combined 1800 kg/$2^\circ$ condition is an interpolation stress cell,
not an exact member of this grid.
Static RACF sets $\Theta_t=\Theta$ and $m_t=\eta q_\alpha$.
The coefficient $\eta$ specifies the operating margin; the analytical
residual-to-safety conversion uses $\eta\geq L_h$.
When the feasible set is empty or the solver fails its feasibility check,
the controller applies bounded maximum braking and records a fallback.
This makes feasibility, intervention, and realized safety independently
observable.

\begin{figure}[t]
\centering
\begin{tikzpicture}[>=Stealth, font=\fontsize{9}{10}\selectfont,
 node distance=5mm, box/.style={draw,align=center,inner sep=4pt,minimum height=9mm}]
 \node[box] (policy) {Frozen policy\\$a_t^\pi=\pi(o_t)$};
 \node[box,right=6mm of policy] (filter) {Robust projection\\$\Theta_t,\,m_t$};
 \node[box,right=6mm of filter] (plant) {ACC\\plant};
 \draw[->] (policy) -- (filter);
 \draw[->] (filter) -- node[above] {$a_t$} (plant);
 \node[box,below=7mm of filter] (cal) {Residual calibration\\and candidate update};
 \draw[->] (cal) -- (filter);
 \draw[->] (plant.south) |- (cal.east);
 \node[right=1mm of plant.south,anchor=north west] {$s_{t+1}$};
\end{tikzpicture}
\caption{RACF combines a frozen policy with calibrated action correction.
The transition observed after execution updates the next control step.}
\label{fig:method}
\end{figure}

\subsection{Adaptive selection and horizon-aware envelopes}
The adaptive variant ranks each candidate by the residual quantile in a
40-transition window, evaluated every ten steps after twenty observations.
It retains at most four candidates with the smallest residual scores.
Observed grade further restricts candidates within $0.75^\circ$ of a grid
value; intermediate grades retain the available grid. A residual trigger
reopens the full candidate set for a bounded recovery period. The threshold
uses a separate 100-transition window with its calibration value as a floor.
All updates consume the transition just completed, so the action at time
$t$ uses information available before its execution (Fig.~\ref{fig:method}).
This empirical adaptation complements the explicit online coverage objectives
studied in adaptive conformal inference \cite{gibbs2021adaptive}.

A horizon-aware alternative calibrates the maximum score of a complete
episode $\tau$:
\begin{equation}
 M_H(\tau)=\max_{0\leq t<H} S(Z_t).
 \label{eq:maxscore}
\end{equation}
Applying~\eqref{eq:quantile} to independent episode maxima yields an
envelope $q_H$. The evaluated fixed-envelope variant instead transfers the
previously selected threshold $B=0.45$ to the 12-policy study. This tests the
effect of an episode-oriented operating margin across policies; the formal
max-score result below describes calibration on complete episodes.
